\paragraph*{Ograniczony $\mu$-operator}
\textit{Definicja z ćwiczeń:} \\
Funkcja \( f : \mathbb{N}^k \to \mathbb{N} \) jest zdefiniowana za pomocą 
\textbf{ograniczonego \(\mu\)-operatora} (\textbf{operacji minimum}) 
z funkcji \( h : \mathbb{N}^k \to \mathbb{N} \) oraz \( g : \mathbb{N}^{k+1} \to \mathbb{N} \), jeśli
\[
    f(\vec{n}) =
    \begin{cases}
        \min \{ m \le h(\vec{n}) : g(\vec{n}, m) = 0 \}, & 
        \text{jeśli } \exists_{m \le h(\vec{n})} \ g(\vec{n}, m) = 0, \\
        0, & \text{wpp.}
    \end{cases}
\]
Zapisujemy wtedy
\[
    f(n_1, \ldots, n_k) = 
\mu \: m_{\le h(n_1, \ldots, n_k)} \ [\, g(n_1, \ldots, n_k, m) = 0 \,].
\]

Na ćwiczeniach (1.5) dowiedliśmy bardzo przydatne twierdzenie:
\begin{theorem}
    Jeśli \( f \) jest zdefiniowana z funkcji \( g, h \in \prec \) za pomocą ograniczonego \(\mu\)-operatora, to \( f \) jest pierwotnie rekurencyjna.
\end{theorem}